\documentclass{article}

\usepackage[margin=1in]{geometry}
\usepackage{amsmath,amsthm,amssymb}
\usepackage[spanish]{babel}
\usepackage{enumitem}

\theoremstyle{definition}
\newtheorem{definition}{Definición}[section]

% Number sets
\newcommand{\R}{\mathbb{R}}
\newcommand{\Z}{\mathbb{Z}}
\newcommand{\N}{\mathbb{N}}
\newcommand{\Q}{\mathbb{Q}}
\newcommand{\C}{\mathbb{C}}

% Topology operators
\newcommand{\cl}[1]{\overline{#1}}              % Closure
\newcommand{\Int}[1]{{#1}^\circ}                % Interior
\newcommand{\bd}{\partial}                       % Boundary
\newcommand{\Ent}[1]{\mathcal{O}(#1)}           % Neighborhoods (Entornos)
\newcommand{\pf}{\mathcal{P}_F}                 % Finite parts
\newcommand{\partes}{\mathcal{P}}               % Power set

% Useful shortcuts
\newcommand{\sii}{\iff}                          % Si y sólo si
\newcommand{\imp}{\implies}                      % Implies
\newcommand{\setdiff}{\smallsetminus}            % Set difference

% Exercise environment
\newenvironment{ejercicio}[1]{\begin{trivlist}
\item[\hskip \labelsep {\bfseries Ejercicio}\hskip \labelsep {\bfseries #1.}]}{\end{trivlist}}

\setlist[enumerate,1]{label=(\alph*), leftmargin=*, itemsep=2pt}

\begin{document}

\title{Topología: Práctica I}
\author{Bustos Jordi\\Teoría de conjuntos}

\maketitle

\begin{ejercicio}{1}
    Sean \(A\) y \(B\) dos conjuntos y \( f:A\rightarrow B\) una función. Sean \( M\subseteq A\) y \(N\subseteq B\). Demostrar:
    \begin{enumerate}
        \item \(M\subseteq f^{-1}(f(M))\). Además, si \(f\) es inyectiva entonces vale la igualdad.
        \item \(f(f^{-1}(N))\subseteq N\). Además, si \(f\) es suryectiva entonces vale la igualdad.
    \end{enumerate}
\end{ejercicio}

\begin{proof}
    (a) Para demostrar que \(M\subseteq f^{-1}(f(M))\), tomamos un elemento \(x \in M\). Por definición de imagen, \(f(x) \in f(M)\). Luego, por definición de preimagen, \(x \in f^{-1}(f(M))\). Por lo tanto, \(M \subseteq f^{-1}(f(M))\).
    Ahora, si \(f\) es inyectiva, queremos demostrar que \(f^{-1}(f(M)) \subseteq M\). Tomamos un elemento \(y \in f^{-1}(f(M))\). Esto significa que \(f(y) \in f(M)\), lo que implica que existe un elemento \(x \in M\) tal que \(f(y) = f(x)\). Dado que \(f\) es inyectiva, concluimos que \(y = x\), y por lo tanto \(y \in M\). Así, \(f^{-1}(f(M)) \subseteq M\).

    (b) Para demostrar que \(f(f^{-1}(N))\subseteq N\), tomamos un elemento \(z \in f(f^{-1}(N))\). Esto significa que existe un elemento \(w \in f^{-1}(N)\) tal que \(f(w) = z\). Por definición de preimagen, \(w \in f^{-1}(N)\) implica que \(f(w) \in N\). Por lo tanto, \(z = f(w) \in N\), y así \(f(f^{-1}(N)) \subseteq N\).
    Ahora, si \(f\) es suryectiva, queremos demostrar que \(N \subseteq f(f^{-1}(N))\). Tomamos un elemento \(n \in N\). Dado que \(f\) es suryectiva, existe un elemento \(a \in A\) tal que \(f(a) = n\). Por definición de preimagen, \(a \in f^{-1}(N)\) porque \(f(a) = n \in N\). Luego, \(f(a) \in f(f^{-1}(N))\), lo que implica que \(n \in f(f^{-1}(N))\). Por lo tanto, \(N \subseteq f(f^{-1}(N))\).
\end{proof}

\begin{ejercicio}{2}
    Sean \(A\) y \(B\) dos conjuntos y \( f:A\rightarrow B\) una función. Sea \( {\{M_{i}\}}_{i\in I} \) una familia no vacía de subconjuntos de \(A\) y \( {\{N_{j}\}}_{j\in J} \) una familia no vacía de subconjuntos de \(B\). ¿Qué relación de inclusión existe entre cada par de conjuntos en los siguientes casos?:
    \begin{enumerate}
        \item \(f(\bigcup_{i\in I}M_{i})\) y \(\bigcup_{i\in I}f(M_{i})\). \(f^{-1}(\bigcup_{j\in J}N_{j})\) y \(\bigcup_{j\in J}f^{-1}(N_{j})\).
        \item \(f(\bigcap_{i\in I}M_{i})\) y \(\bigcap_{i\in I}f(M_{i})\). \(f^{-1}(\bigcap_{j\in J}N_{j})\) y \(\bigcap_{j\in J}f^{-1}(N_{j})\).
    \end{enumerate}
\end{ejercicio}

\begin{proof}

\end{proof}

\begin{ejercicio}{3}
    Sean \(A\) y \(B\) dos conjuntos, \(N\subseteq B\) y \(f:A\rightarrow B\) una función. ¿Qué relación de inclusión existe entre los siguientes conjuntos?:
    \(f^{-1}(B-N)\) y \(A-f^{-1}(N)\).
\end{ejercicio}

\begin{proof}
    \( x \in f^{-1}(B-N) \iff f(x) \in B-N \iff f(x) \notin N \iff x \notin f^{-1}(N) \iff x \in A - f^{-1}(N) \).
\end{proof}

\begin{ejercicio}{4}
    Sean \(A\) y \(B\) dos conjuntos, \(M\subseteq A\) y \(f:A\rightarrow B\) una función. ¿Qué relación de inclusión existe entre los siguientes conjuntos?:
    \(f(A-M)\) y \(B-f(M)\).
    ¿Qué propiedades sobre \(f\) son necesarias para que sea válida la igualdad de los conjuntos anteriores?.
\end{ejercicio}

\begin{proof}
    Para ver que \( f(A-M) \subseteq B-f(M) \), tomamos un elemento \( y \in f(A - M) \subseteq B \), luego existe un elemento \( x \in A - M \) tal que \( f(x) = y \). Dado que \( x \notin M \), no es posible que \( f(x) \in f(M) \), por lo que \( y \in B - f(M) \).
    Para ver que la otra contención no vale consideremos el siguiente contraejemplo, sea \( A = \{ a, b \} \), \( B = \{ c, d \} \), \( M = \{ a \} \) y \( f : A \to B \) definida por \( f(a) = c \) y \( f(b) = c \). Entonces \( f(A - M) = f(\{ b \}) = \{ c \} \) y \( B - f(M) = B - f(\{ a \}) = B - \{ c \} = \{ d \} \), por lo que \( f(A - M) \not\subseteq B - f(M) \).
    Para que se cumpla la igualdad \( f(A - M) = B - f(M) \) es necesario que \( f \) sea inyectiva. En efecto, si \( f \) es inyectiva, entonces \( f(M) \) es un subconjunto de \( B \) que no contiene a ningún elemento de \( f(A - M) \), lo que implica que \( f(A - M) \subseteq B - f(M) \). Por otro lado, si \( y \in B - f(M) \), entonces no existe ningún elemento \( x \in M \) tal que \( f(x) = y \). Dado que \( f \) es inyectiva, esto implica que no existe ningún elemento \( x \in A - M \) tal que \( f(x) = y \), lo que implica que \( y \in f(A - M) \). Por lo tanto, \( B - f(M) \subseteq f(A - M) \). Combinando ambas contenciones, concluimos que \( f(A - M) = B - f(M) \) si \( f \) es inyectiva.
\end{proof}

\begin{ejercicio}{5}
    Para \(f:A\rightarrow B\), \(g:B\rightarrow C\) y \(M\subseteq C\). Entonces \( {(g\circ f)}^{-1}(M)=f^{-1}(g^{-1}(M)) \).
\end{ejercicio}

\begin{proof}
    Es claro que \( {(g \circ f)}^{-1} \) es la inversa de \( g \circ f \), luego \[
        f^{-1}(g^{-1}) \circ (g \circ f) = id_A = (g \circ f) \circ f^{-1}(g^{-1}) \implies f^{-1}(g^{-1}) = {(g \circ f)}^{-1}
    \]
    por la unicidad de la inversa de una función.
\end{proof}

\begin{ejercicio}{6}
    Sea \( X \) un conjunto y \( F:\mathcal{P}(X)\rightarrow\mathcal{P}(X)\) tal que si \(U\subseteq V\subseteq X\) entonces \(F(U)\subseteq F(V)\). Demostrar que existe \(Z\in\mathcal{P}(X)\) tal que \(F(Z)=Z\).
\end{ejercicio}

\begin{proof}
    Definamos la siguiente familia de conjuntos \[
        \mathcal{F} = \{ A \subseteq X : A \subseteq F(A) \}
    \]
    Nuestro candidato es \[
        Z = \bigcup_{A \in \mathcal{F}} A
    \]
    En efecto, veamos que \( Z \subseteq F(Z) \). Dado \( A \in Z \) existe un \( A_0 \in \mathcal{F} \) tal que \( A \in A_0 \). Por definición de \( \mathcal{F} \) se tiene que \( A_0 \subseteq F(A_0) \) y
    como \( A_0 \) es parte de la unión que conforma a \( Z \) se tiene que \( A_0 \subseteq Z \) y tenemos que \( F(A_0) \subseteq F(Z) \) por hipotesis, luego \[
        A \in A_0 \subseteq F(A_0) \subseteq F(Z) \implies Z \subseteq F(Z)
    \]
    y, por lo tanto, \( Z \subseteq F(Z) \).
    Por otro lado, para ver que \( F(Z) \subseteq Z \) basta notar que \( F(Z) \subseteq F(F(Z)) \) y esto implica que \( F(Z) \in \mathcal{F} \) y por lo tanto \( F(Z) \subseteq Z \).
\end{proof}

\begin{ejercicio}{7}
    Demostrar la equivalencia entre:
    \begin{itemize}
        \item AdE
        \item AdE 2
        \item Para cualquier familia \( {\{X_{i}\}}_{i\in I}\) de conjuntos no vacios y disjuntos dos a dos existe un conjunto \(Y\) que cumple \(|Y\cap X_{i}|=1\) para cada \(i\in I\).
    \end{itemize}
\end{ejercicio}

\begin{definition}[AdE]
    Sea \( A \neq \varnothing \) un conjunto. Entonces existe una función \[
        e: \mathcal{P}_0(A) \to A \quad \text{tal que} \quad e(X) \in X \text{ para todo } X \in \mathcal{P}_0(A)
    \]
\end{definition}

\begin{definition}[AdE 2]
    Dada una familia de conjuntos \( {\{ X_i \}}_{i \in I} \) vale que \[
        \text{si } X_j \neq \varnothing \quad \forall j \in I \implies \prod_{i \in J} X_i \neq \varnothing
    \]
\end{definition}

\begin{proof}
    Para demostrar que AdE implica AdE 2, supongamos que AdE es cierta. Luego, dada la familia \( {\{ X_i \}}_{i \in I} \) con \( X_j \neq \varnothing \) para todo \( j \in I \), sea \( e \) la función de elección del axioma:
    definimos \[
        f: I \to \bigcup_{i \in I} X_i \quad \text{tal que} \quad f(i) = e(X_i) \in X_i
    \]
    Esto implica que \( f \in \prod_{i \in I} X_i \), por lo tanto \( \prod_{i \in I} X_i \neq \varnothing \).

    Veamos que AdE 2 implica la tercera afirmación. Dada la familia \( {\{ X_i \}}_{i \in I} \) de conjuntos no vacíos y disjuntos dos a dos, por AdE 2 existe una función \( f \in \prod_{i \in I} X_i \).
    Definamos ahora el conjunto \( Y = \{ f(i) : i \in I \} \) que es la imagen de \( f \). Dado que \( f(i) \in X_i \) para cada \( i \in I \) y que \( f(i) \neq f(j) \) si \( i \neq j \) por ser disjuntos dos a dos,
    se tiene que \( |Y \cap X_i| = 1 \) para cada \( i \in I \).

    Finalmente veamos que la tercera afirmación implica AdE. Dado \( \mathcal{P_0}(A) = \mathcal{F} \) definamos la familia \[
        \mathcal{\bar{F}} = \{ X \times \{ X \} : X \in \mathcal{F} \}
    \]
    Notar que la familia \( \mathcal{\bar{F}} \) es una familia de conjuntos no vacíos y disjuntos dos a dos. Por hipotesis existe un conjunto \( Y \) tal que \( |Y \cap (X \times \{ X \})| = 1 \).
    Digamos que \( Y \cap (X \times \{ X \}) = (x, X) \) para cada \( X \in \mathcal{F} \). Definamos ahora la función de elección \( e : \mathcal{F} \to A \) por \( e(X) = x \). Es claro que \( e(X) \in X \) para cada \( X \in \mathcal{F} \), por lo tanto AdE se cumple.
\end{proof}

\begin{ejercicio}{8}
    Dar contraejemplos para mostrar que no son válidas las recíprocas de las siguientes implicaciones:
    Buen orden \( \implies \) Orden total \( \implies \) Lattice \( \implies \) Orden dirigido (superiormente e inferiormente)
\end{ejercicio}

\begin{definition}[Buen orden]
    Un orden \( (A, \leq) \) es un buen orden si todo subconjunto no vacío de \( A \) tiene un primer elemento.
\end{definition}

\begin{definition}[Orden total]
    Un orden \( (A, \leq) \) es un orden total si para todo \( a, b \in A \) se cumple que \( a \leq b \) o \( b \leq a \).
\end{definition}

\begin{definition}[Lattice]
    Un orden \( (A, \leq) \) es un lattice si para todo \( a, b \in A \) existen el supremo \( a \vee b \) y el ínfimo \( a \wedge b \).
\end{definition}

\begin{definition}[Orden Dirigido]
    Un orden \( (A, \leq) \) es un orden dirigido superiormente si para todo \( a, b \in A \) existe \( c \in A \) tal que \( a \leq c \) y \( b \leq c \). De manera similar, es un orden dirigido inferiormente si para todo \( a, b \in A \) existe \( d \in A \) tal que \( d \leq a \) y \( d \leq b \).
\end{definition}

\begin{proof}
    Para el ejemplo de un orden total que no sea un buen orden consideremos \( (\Z, \leq) \), el orden es total pero el subconjunto \( \N^- \) no tiene primer elemento, por lo tanto no es un buen orden. \\
    Un Lattice que no es Orden total es el siguiente, dado \( X \) un conjunto con al menos dos elementos, \( (\mathcal{P}(X), \subseteq) \) pues los singueletes son incomparables. \\
    Para el contraejemplo de un orden dirigido que no es un lattice consideremos el siguiente conjunto \( A = \{ a, b, c, d, e \} \) con el orden dado por \begin{itemize}
        \item \( a \leq c \) y \( b \leq c \)
        \item \( a \leq d \) y \( b \leq d \)
        \item \( c \leq e \) y \( d \leq e \)
        \item Las cotas superiores de \( \{ a, b \} \) son \( \{ c, d, e \} \) y no existe una cota superior mínima pues \( c \) y \( d \) son incomparables, por lo tanto no es un lattice.
    \end{itemize}
\end{proof}

\begin{ejercicio}{9}
    Demostrar que el orden definido entre los números cardinales es un orden total.
\end{ejercicio}

\begin{proof}
    Sean A y B conjuntos cualesquiera quiero ver que existe \( f : A \to B \) inyectiva o \( g: A \to B \) suryectiva para poder garantizar que \( |A| \leq |B| \) o \( |B| \leq |A| \).

    Para esto consideremos: \[
        F = \{  (X, f) : X \subseteq A \text{ y } f: X \to B \text{ inyectiva} \}
    \]
    Ahora dotemos a \( F \) con el orden \( (X, f) \leq ( Y, g ) \) si \( X \subseteq Y \) y \( g(x) = f(x) \) para todo \( x \in X \), es decir \( g \) extiende a \( f \).
    Para aplicar el LdZ necesitamos ver que toda cadena tiene un elemento maximal, pero esto es sencillo, dada una cadena podemos tomar la unión de todos los dominios \( U \) de la misma y definir la función \( h: U \to B \) inyectiva tomando el valor que le daba cualquiera de las funciones en la cadena.
    Luego, es claro que \( (U, h) \) es cota superior de la cadena y esto nos dice que existe \( (M, f^*) \) elemento maximal de \( F \) con \( M \subseteq A \).

    Consideremos los dos casos posibles, si \( M = A \) entonces \( f^* : M = A \to B \) es inyectiva, por otro lado si \( M \neq A \) entonces existe al menos un elemento \( a \in A - M \). Si existiese \( b \in B \) tal que \( f^*(a) = b \) entonces \( (M, f^*) \) no sería maximal.
    Esto nos dice que la imagen de \( f^* \) debe ser todo \( B \), entonces podemos concluir que \( |A| \leq |B| \) o \( |B| \leq |A| \) y por lo tanto el orden definido entre los números cardinales es un orden total.
\end{proof}

\begin{ejercicio}{10}
    Demostrar las siguientes propiedades referida a conjuntos numerables.
    \begin{enumerate}
        \item El producto cartesiano de \(n\) conjuntos numerables es también un conjunto numerable.
        \item La unión de una cantidad numerable de conjuntos numerables es un conjunto numerable.
        \item Todo subconjunto de un conjunto numerable es un conjunto numerable.
        \item Partes finitas de un conjunto numerable es también un conjunto numerable.
    \end{enumerate}
\end{ejercicio}

\begin{proof}
    (a) Se demuestra por inducción sobre la cantidad de conjuntos \(n\).

    \textbf{Caso base (\(n=2\)):} Sean \(A\) y \(B\) dos conjuntos numerables. Existen funciones inyectivas \(f: A \to \N\) y \(g: B \to \N\). Definimos una nueva función para el producto cartesiano \(h: A \times B \to \N\) de la siguiente manera:
    \(h(a, b) = 2^{f(a)} \cdot 3^{g(b)}\)

    Por el Teorema Fundamental de la Aritmética (la factorización en números primos es única), si \(2^{f(a)} \cdot 3^{g(b)} = 2^{f(a')} \cdot 3^{g(b')}\), entonces obligatoriamente \(f(a) = f(a')\) y \(g(b) = g(b')\). Como \(f\) y \(g\) son inyectivas, esto implica que \(a = a'\) y \(b = b'\). Por lo tanto, \(h\) es inyectiva y \(A \times B\) es numerable.

    \textbf{Paso inductivo:} Asumimos que el producto cartesiano de \(k\) conjuntos numerables es numerable (hipótesis inductiva). Queremos probarlo para \(k+1\) conjuntos. Sabemos que:
    \(\prod_{i=1}^{k+1} A_i \cong \left( \prod_{i=1}^{k} A_i \right) \times A_{k+1}\)

    Por hipótesis inductiva, el término \(\prod_{i=1}^{k} A_i\) es un conjunto numerable. \(A_{k+1}\) también lo es por definición. Al agruparlos así, tenemos el producto cartesiano de \textit{dos} conjuntos numerables, que ya demostramos en el caso base que es numerable. Por el principio de inducción, la propiedad se cumple para todo \(n\).

    (b) Sea \( {\{A_i \}}_{i \in \N}\) una familia numerable de conjuntos numerables. Queremos demostrar que su unión \(U = \bigcup_{i \in \N} A_i\) es numerable.

    Como cada \(A_i\) es numerable, podemos listar sus elementos. Sea \(a_{i,j}\) el \(j\)-ésimo elemento del conjunto \(i\)-ésimo.
    Para definir una función inyectiva \(F: U \to \N\), tomamos un elemento cualquiera \(x \in U\). Como \(x\) puede pertenecer a varios conjuntos de la familia, elegimos el menor índice \(i\) tal que \(x \in A_i\), y luego elegimos el índice \(j\) correspondiente a su posición en ese conjunto (es decir, \(x = a_{i,j}\)).

    Definimos la función:
    \(F(x) = 2^i \cdot 3^j\)

    Dado que para cada \(x\) los valores de \(i\) y \(j\) elegidos son únicos, y la factorización en primos es única, \(F(x) = F(y)\) implica que tienen los mismos índices \(i\) y \(j\), por lo tanto \(x = y\). Así, \(F\) es inyectiva y la unión \(U\) es numerable.

    (c) Sea \(A\) un conjunto numerable y sea \(B \subseteq A\) un subconjunto cualquiera de \(A\).

    Dado que \(A\) es numerable, existe una función inyectiva \(f: A \to \N\).
    Consideramos la restricción de \(f\) al subconjunto \(B\), denotada como \(f|_B : B \to \N\). Esta función simplemente toma un elemento \(x \in B\) y le aplica la misma regla \(f(x)\).

    Si tomamos \(x, y \in B\) tales que \(f|_B(x) = f|_B(y)\), esto significa que \(f(x) = f(y)\). Como sabemos que \(f\) original es inyectiva en todo el dominio \(A\), concluimos que \(x = y\).
    Al haber encontrado una función inyectiva de \(B\) a \(\N\), queda demostrado que \(B\) es numerable.

    (d) Sea \(A\) un conjunto numerable. Queremos demostrar que el conjunto de todas sus partes finitas, \(\mathcal{P}_{fin}(A)\), es numerable.

    Podemos clasificar los subconjuntos finitos por su tamaño. Sea \(S_k\) el conjunto de todos los subconjuntos de \(A\) que tienen exactamente \(k\) elementos.
    Cualquier subconjunto de tamaño \(k\) puede formarse eligiendo \(k\) elementos de \(A\), lo cual puede representarse ordenando arbitrariamente sus elementos como una \(k\)-tupla, es decir, un elemento de \(A^k\).

    Dado que \(A\) es numerable, \(A^k\) es el producto cartesiano de \(k\) conjuntos numerables, el cual sabemos que es numerable por la propiedad (a). Como todo elemento de \(S_k\) proviene de \(A^k\), \(S_k\) debe ser numerable (es la imagen de un conjunto numerable).

    Finalmente, el conjunto de todas las partes finitas es la unión de todos estos conjuntos según su tamaño:
    \(\mathcal{P}_{fin}(A) = \bigcup_{k=0}^{\infty} S_k\)

    Esta es la unión numerable (indexada por \(k \in \N_0\)) de conjuntos numerables (\(S_k\)). Por la propiedad (b), esta unión es un conjunto numerable.

\end{proof}

\begin{ejercicio}{11}
    Demostrar que \( {\{0,1\}}^{\N} \) no es numerable.
\end{ejercicio}

\begin{proof}
    Procederemos por reducción al absurdo utilizando el argumento diagonal de Cantor. Supongamos que el conjunto \( S = {\{0,1\}}^\N \) es numerable.

    Esto implica que existe una biyección entre los números naturales \( \N \) y \( S \). Por lo tanto, podemos enumerar todas las sucesiones de \(S\) como \(s_1, s_2, s_3, \dots \), donde cada sucesión \(s_n\) se puede escribir en función de sus componentes como:
    \[s_n = (s_{n,1}, s_{n,2}, s_{n,3}, \dots)\]
    con \( s_{n,i} \in \{0,1\} \) para todo \(n, i \in \N\).

    Construyamos ahora una nueva sucesión \(d = (d_1, d_2, d_3, \dots)\) tal que su \(n\)-ésimo término se define invirtiendo el \(n\)-ésimo término de la \(n\)-ésima sucesión de nuestra lista. Formalmente:
    \[d_n = 1 - s_{n,n}\]
    (es decir, si \(s_{n,n} = 0\), entonces \(d_n = 1\); y si \(s_{n,n} = 1\), entonces \(d_n = 0\)).

    Dado que cada \(d_n \in \{0,1\} \), es claro que la sucesión \(d \in S\). Sin embargo, por construcción, la sucesión \(d\) difiere de cada sucesión \(s_n\) en al menos el \(n\)-ésimo término (ya que \(d_n \neq s_{n,n}\)).

    Por lo tanto, \( d \neq s_n \) para todo \(n \in \N \), lo que significa que \(d\) no está en nuestra enumeración original. Esto contradice el supuesto de que habíamos enumerado \emph{todas} las sucesiones de \(S\). Esta contradicción demuestra que nuestro supuesto inicial era falso, por lo que concluimos que \( {\{0,1\}}^\N \) no es numerable.
\end{proof}

\begin{ejercicio}{12}
    Demostrar que si \(A\) es infinito y \(F\) es finito entonces \(|A\cup F|=|A|\).
\end{ejercicio}

\begin{proof}
    Por el teorema de Cantor-Bernstein-Schröder, para demostrar que \(|A\cup F|=|A|\) es suficiente mostrar que \(A\cup F\) se inyecta en \(A\) y que \(A\) se inyecta en \(A\cup F\).
    Consideremos por un lado \( f : A \to A \cup F \) definida por \( f(a) = a \) para todo \( a \in A \). Esta función es claramente inyectiva.

    Veamos ahora como construir una función inyectiva \( g : A \cup F \to A \). Dado que \( F \) es finito podemos considerar \[
        F = \{ x_1 , \ldots , x_n \}
    \]
    Por el \textbf{Teorema 1.6.1} del apunte sabemos que \( \exists B \subset A \) tal que \( |B| = \aleph_0 \), es decir, \(B\) es un subconjunto infinito numerable de \(A\). Por lo tanto, podemos enumerar los elementos de \(B\) como \( B = \{ b_1 , b_2 , b_3 , \ldots \} \).
    Definamos ahora la función \( g : A \cup F \to A \) de la siguiente manera:
    \[
        g(x) = \begin{cases}
            x       & \text{si } x \in A - B                                   \\
            b_i     & \text{si } x = x_i \in F \text{ para } i = 1, \ldots , n \\
            b_{n+i} & \text{si } x = b_i \in B \text{ para } i = 1, 2, \ldots
        \end{cases}
    \]
    Es fácil verificar que \( g \) es inyectiva. Luego por el teorema de Cantor-Bernstein-Schröder, concluimos que \( |A\cup F|=|A| \).
\end{proof}

\section*{Ejercicios complementarios}

\begin{ejercicio}{1}
    Demostrar que todo morfismo \( f:A\rightarrow A\) (\(A \) un retículo completo) tiene al menos un punto fijo.
\end{ejercicio}

\begin{ejercicio}{2}
    Demostrar que existe al menos un conjunto de números reales que no es medible Lebesgue.
\end{ejercicio}

\begin{ejercicio}{3}
    Demostrar que en todo conjunto no vacío puede definirse una estructura de grupo.
\end{ejercicio}

\begin{ejercicio}{4}
    Demostrar que el cardinal de un conjunto infinito es igual al cardinal de su conjunto partes finitas.
\end{ejercicio}

\end{document}
